\section{Introduction}

Decentralized sequential decision-making serves as the computational backbone for modern complex systems, ranging from high-frequency automated auctions and smart-grid management to autonomous fleet coordination \cite{shoham2008multiagent, sutton2018reinforcement}. In these environments, independent agents must navigate high-dimensional state spaces while simultaneously adapting to the evolving strategies of competitors. Despite the success of Multi-Agent Reinforcement Learning (MARL) in controlled settings \cite{lowe2017multi}, deployment in real-world, high-frequency competitive domains remains hazardous. The primary challenge arises from the inherent non-stationarity of the environment; as agents learn, the transition dynamics perceived by any single agent shift, leading to ``fragile equilibria'' that collapse under minor distributional perturbations or strategic shifts \cite{hernandez2019survey, papoudakis2021agent}.

Current paradigms in MARL are predominantly optimization-centric, treating the objective as a black-box maximization of cumulative reward or specific engineering metrics like Return on Ad Spend (ROAS). While effective for local performance tuning, this approach ignores the underlying structural dependencies that govern system behavior. Such models often succumb to the ``correlation trap,'' where agents overfit to spurious patterns in historical trajectories that lack predictive power when the system moves out-of-distribution (OOD) \cite{pitis2022counterfactual, arjovsky2019invariant}. In competitive landscapes, this leads to a ``Latency Illusion,'' where agents optimize for immediate feedback loops while remaining blind to the latent causal mechanisms that drive long-term equilibrium stability. We argue that the lack of a formal causal foundation is the fundamental bottleneck preventing the realization of robust, autonomous decentralized intelligence.

To bridge this gap, we move beyond the engineering focus on metric optimization toward a scientific framework grounded in causal discovery. We introduce \textbf{Causal Multi-Agent Reinforcement Learning (Causal MARL)}, a theoretical framework that defines agent interactions through the lens of structural causal models (SCMs) \cite{pearl2009causality}. Unlike standard MARL, which learns policies $\pi(a|s)$ based on conditional correlations, Causal MARL seeks to identify \textit{invariant causal relationships}—mechanisms that remain stable even when other agents change their strategies or the environment undergoes external interventions \cite{peters2016causal, scholkopf2021causality}. This shift allows agents to distinguish between transient environmental noise and the persistent causal levers that influence system-wide equilibria.

Our vision is summarized in Figure 1 [Note: See Figure 1 for the conceptual flow], which illustrates the transition from correlational MARL to causally-aware equilibrium seeking. By discovering the latent causal graph of the decision-making process, agents can anticipate the downstream effects of their actions on the collective state, leading to what we term \textbf{Equilibrium Robustness}. This property ensures that the resulting Nash or Stackelberg equilibria are not merely artifacts of a specific training distribution but are anchored in the structural invariants of the system.

Specifically, this paper makes three primary contributions:
\begin{enumerate}
    \item \textbf{Theoretical Foundation}: We formally define Causal MARL and provide necessary conditions for discovering invariant causal structures in non-ergodic, multi-agent trajectories.
    \item \textbf{Invariant Causal Discovery (ICD) Algorithm}: We propose a novel class of algorithms that utilize intervention-based reasoning to prune spurious correlations and identify robust causal drivers in decentralized systems.
    \item \textbf{Equilibrium Stability Guarantees}: We provide theoretical and empirical evidence that agents equipped with causal awareness achieve significantly higher stability and predictability in competitive settings compared to SOTA MARL baselines \cite{yu2022surprising, gao2024causal}.
\end{enumerate}

By centering causality in the multi-agent loop, we offer a path toward decentralized systems that are not only efficient but fundamentally reliable. The following sections detail our theoretical framework, the algorithmic implementation of Causal MARL, and extensive empirical evaluations across complex, non-stationary environments.